%============================================================
%  ch10_DF.tex
%  Chapter 10: The Dirac--Frenkel Variational Principle
%
%  Tensor Formats in Banach Spaces
%  Antonio Falcó, CEU Universities
%
%  Based on: FHN2019 (Sec. 5), FHN2023 (Sec. 5)
%  Estimated length: ~35 pp.
%============================================================

\chapter{The Dirac--Frenkel Variational Principle}
\label{chap:DF}

The geometric framework developed in Chapters~\ref{chap:tucker_manifold}
and~\ref{chap:TB_manifold} has a direct dynamical application: it
provides the rigorous foundation for \emph{model reduction} of
differential equations on tensor Banach spaces. Given an abstract
Cauchy problem on $\VDnorm$, one looks for an approximate solution
that evolves on the manifold $\Tucker{\fr}{\VD}$ (or $\FTv$).
The natural way to constrain the dynamics to the manifold is the
\emph{Dirac--Frenkel variational principle}: at each time $t$, the
velocity $\dot{\bv}_r(t)$ is chosen as the best approximation of the
vector field $\bF(t, \bv_r(t))$ in the tangent space $\ZD{\bv_r(t)}$.

This principle was originally proposed by Dirac and Frenkel in 1930 for
the time-dependent Schrödinger equation in Hilbert spaces. The present
chapter extends it to general Banach tensor spaces, establishing the
existence of the reduced ODE and characterising its solutions.
Sections~\ref{sec:DF:cauchy}--\ref{sec:DF:hilbert} set up the
abstract framework. Section~\ref{sec:DF:tucker} proves the Tucker
version (Theorems~5.2 and~5.4 of~\cite{FHN2019}).
Section~\ref{sec:DF:treebased} extends to the tree-based setting
(\cite{FHN2023}).

%------------------------------------------------------------
\section{The abstract Cauchy problem on tensor Banach spaces}
\label{sec:DF:cauchy}
%------------------------------------------------------------

Let $(V_\alpha, \|\cdot\|_\alpha)$ be normed spaces for $\alpha \in D$,
and let $\VDnorm = {}_{\|\cdot\|_D}\bigotimes_{\alpha \in D} V_\alpha$
be a reflexive Banach tensor space satisfying condition~\condINJ.
Consider the abstract Cauchy problem:
\begin{align}
    \dot{\bu}(t) &= \bF(t, \bu(t)), \quad t \geq 0,
    \label{eq:DF:ode} \\
    \bu(0) &= \bu_0,
    \label{eq:DF:ic}
\end{align}
where $\bu_0 \neq \mathbf{0}$, and $\bF: [0, \infty) \times \VDnorm
\to \VDnorm$ satisfies the usual conditions for existence and uniqueness
of solutions (e.g., Lipschitz continuity in the second argument).

The goal is to approximate $\bu(t)$ by a \emph{differentiable curve}
$t \mapsto \bv_r(t)$ lying in $\Tucker{\fr}{\VD}$ for some fixed
$\fr \in \AD{\VD}$ with $\fr \neq \mathbf{0}$, with initial condition
$\bv_r(0) = \bv_0 \in \Tucker{\fr}{\VD}$ chosen as an approximation
of $\bu_0$.

\begin{remark}
\label{rem:DF:initial}
By Theorem~\ref{thm:best_approx_tucker}, a best approximation
$\bv_0 \in \Tucker{\leq\fr}{\VD}$ of $\bu_0$ exists when $\VDnorm$
is reflexive. The initial condition $\bv_0 \in \Tucker{\fr}{\VD}$
(with exact rank $\fr$) can always be arranged by taking
$\bv_0 \in \Tucker{\fr'}{\VD}$ for some $\fr' \leq \fr$ and
adjusting the rank.
\end{remark}

%------------------------------------------------------------
\section{The principle in Hilbert spaces}
\label{sec:DF:hilbert}
%------------------------------------------------------------

When $\VDnorm$ is a Hilbert space with inner product
$\langle \cdot, \cdot \rangle_D$, the constraint to $\Tucker{\fr}{\VD}$
takes the classical form. Since $\ZD{\bv_r(t)} \in \GrassB{\VDnorm}$
(Theorem~\ref{thm:tucker_immersion}), the orthogonal projection
$\Pmetric{\bv_r(t)}$ onto $\ZD{\bv_r(t)}$ is well-defined and the
\emph{Dirac--Frenkel equation} reads:
\begin{equation}
\label{eq:DF:hilbert}
    \dot{\bv}_r(t) = \Pmetric{\bv_r(t)}\bigl(\bF(t, \bv_r(t))\bigr),
\end{equation}
which is equivalent to the variational condition:
\begin{equation}
\label{eq:DF:variational}
    \langle \dot{\bv}_r(t) - \bF(t, \bv_r(t)),\, \dot{\bv} \rangle_D = 0
    \quad \text{for all } \dot{\bv} \in \ZD{\bv_r(t)}.
\end{equation}
The MCTDH (Multi-Configuration Time-Dependent Hartree) method of
quantum dynamics~\cite{Lubich2008} is precisely~\eqref{eq:DF:hilbert}
with $\VDnorm = L^2(\RR^{3d})$ and $\Tucker{\fr}{\VD}$ the Hartree
manifold ($\fr = (1,\ldots,1)$).

%------------------------------------------------------------
\section{Projections in Banach spaces}
\label{sec:DF:projection}
%------------------------------------------------------------

In a general Banach space, there is no canonical inner product and hence
no canonical projection. Two notions replace the orthogonal projection.

\subsection*{Metric projection}

\begin{definition}
\label{def:metric_projection}
For $\bv_r(t) \in \Tucker{\fr}{\VD}$ and $t \in I$, the
\emph{metric projection} $\Pmetric{\bv_r(t)}$ maps
$\dot{\bu} \in \VDnorm$ to the nearest point in $\ZD{\bv_r(t)}$:
\begin{equation*}
    \Pmetric{\bv_r(t)}(\dot{\bu})
    \coloneqq \argmin_{\dot{\bv} \in \ZD{\bv_r(t)}}
    \|\dot{\bu} - \dot{\bv}\|_D.
\end{equation*}
\end{definition}

By Corollary~\ref{cor:best_proj_tangent}, when $\VDnorm$ is reflexive
the metric projection exists. When $\VDnorm$ is strictly convex it is
unique (since $\ZD{\bv_r(t)}$ is convex). The metric projection is,
in general, a \emph{nonlinear} operator on Banach spaces.

\subsection*{Normalised duality mapping}

The duality mapping $J: \VDnorm \to 2^{\dual{\VDnorm}}$ is defined by
\begin{equation*}
    J(\bu) \coloneqq \bigl\{f \in \dual{\VDnorm} :
    \langle \bu, f \rangle = \|\bu\|_D^2 = \|f\|_*^2\bigr\}.
\end{equation*}
It is single-valued when $\VDnorm$ is smooth (Definition~\ref{def:smooth}),
and uniformly continuous on bounded sets when $\VDnorm$ is uniformly smooth.
In a Hilbert space, $J$ coincides with the Riesz isomorphism.

\begin{theorem}[Metric projection characterisation,
{\cite[Theorem~5.2]{FHN2019}}]
\label{thm:metric_projection}
Let $\VDnorm$ be reflexive and strictly convex, and assume \condINJ.
For $\bv_r(t) \in \Tucker{\fr}{\VD}$ and fixed $t \in I$, the
following are equivalent:
\begin{enumerate}
\item[\textup{(a)}] $\dot{\bv}_r(t) = \Pmetric{\bv_r(t)}
    (\bF(t,\bv_r(t)))$.
\item[\textup{(b)}] $\langle \dot{\bv}_r(t) - \dot{\bv},\,
    J(\bF(t,\bv_r(t)) - \dot{\bv}_r(t)) \rangle \geq 0$
    for all $\dot{\bv} \in \ZD{\bv_r(t)}$.
\item[\textup{(c)}] $\langle \dot{\bv},\,
    J(\bF(t,\bv_r(t)) - \dot{\bv}_r(t)) \rangle = 0$
    for all $\dot{\bv} \in \ZD{\bv_r(t)}$.
\end{enumerate}
\end{theorem}

\begin{proof}
This is a direct consequence of the classical characterisation of the
metric projection onto a closed convex set in a reflexive, strictly
convex Banach space, combined with~\cite[Proposition~2.10]{Alber1996}.
Condition~(c) recovers the Hilbert-space variational form: when $J$
is the Riesz map, (c) becomes the orthogonality condition~\eqref{eq:DF:variational}.
\end{proof}

\subsection*{Generalised projection}

When $\VDnorm$ is additionally smooth, an alternative projection is
available. Define $\phi: \VDnorm \times \VDnorm \to \RR$ by
\begin{equation*}
    \phi(\bu, \bv) = \|\bu\|_D^2 - 2\langle \bu, J(\bv) \rangle + \|\bv\|_D^2.
\end{equation*}
The \emph{generalised projection} $\Pi_{\bv_r(t)}$ minimises $\phi$
over $\ZD{\bv_r(t)}$:
\begin{equation*}
    \Pi_{\bv_r(t)}(\dot{\bu})
    \coloneqq \argmin_{\dot{\bv} \in \ZD{\bv_r(t)}}
    \phi(\dot{\bv}, \dot{\bu}).
\end{equation*}
It coincides with the metric projection when $\VDnorm$ is a Hilbert
space.

%------------------------------------------------------------
\section{The Dirac--Frenkel principle for the Tucker manifold}
\label{sec:DF:tucker}
%------------------------------------------------------------

\begin{theorem}[Dirac--Frenkel for Tucker, {\cite[Theorems~5.2,~5.4]{FHN2019}}]
\label{thm:DF_tucker}
Let $\VDnorm$ be reflexive and strictly convex, and assume \condINJ.
Let $\fr \in \AD{\VD}$ with $\fr \neq \mathbf{0}$, and let
$\bv_0 \in \Tucker{\fr}{\VD}$.
\begin{enumerate}
\item[\textup{(Existence)}] There exists a differentiable curve
    $t \mapsto \bv_r(t)$ in $\Tucker{\fr}{\VD}$, defined for $t$ in a
    neighbourhood of $0$, satisfying
    \begin{equation}
    \label{eq:DF_tucker}
        \dot{\bv}_r(t) = \Pmetric{\bv_r(t)}\bigl(\bF(t, \bv_r(t))\bigr),
        \qquad \bv_r(0) = \bv_0.
    \end{equation}

\item[\textup{(Characterisation)}] The curve $t \mapsto \bv_r(t)$
    satisfies~\eqref{eq:DF_tucker} if and only if $\dot{\bv}_r(t)
    \in \ZD{\bv_r(t)}$ and
    \begin{equation}
    \label{eq:DF_tucker_char}
        \|\dot{\bv}_r(t) - \bF(t, \bv_r(t))\|_D
        = \min_{\dot{\bv} \in \ZD{\bv_r(t)}}
        \|\dot{\bv} - \bF(t, \bv_r(t))\|_D.
    \end{equation}
\end{enumerate}
\end{theorem}

\begin{proof}[Sketch]
\textbf{Existence.} The reduced ODE~\eqref{eq:DF_tucker} is an ODE on
the Banach manifold $\Tucker{\fr}{\VD}$. By
Theorem~\ref{thm:tucker_immersion}, $\ZD{\bv_r(t)}$ is a split
subspace of $\VDnorm$ for each $\bv_r(t) \in \Tucker{\fr}{\VD}$.
The right-hand side of~\eqref{eq:DF_tucker} depends smoothly on
$\bv_r(t)$ (the metric projection onto a smoothly-varying split
subspace is smooth), so local existence and uniqueness follow from the
ODE existence theorem on Banach manifolds~\cite[Theorem~4.1.5]{Lang1999}.

\textbf{Characterisation.} Since $\ZD{\bv_r(t)}$ is closed and
convex, Corollary~\ref{cor:best_proj_tangent} gives existence of the
minimum in~\eqref{eq:DF_tucker_char}. The equivalence with the metric
projection condition~\eqref{eq:DF_tucker} follows from the uniqueness
(strict convexity) and Theorem~\ref{thm:metric_projection}.
\end{proof}

\begin{remark}[Kinematic equations]
\label{rem:kinematic}
In local chart coordinates $(L_\alpha, C^{(D)})$ at $\bv_r(t)$,
the reduced ODE~\eqref{eq:DF_tucker} takes the form of a system
of coupled equations for the subspace parameters $L_\alpha(t)$ and
the core tensor $C^{(D)}(t)$. The tangent formula~\eqref{eq:tangent_formula}
makes this explicit:
\begin{equation*}
    \dot{\bv}_r(t)
    = \sum_{(i_\alpha)} \dot{C}^{(D)}_{(i_\alpha)} \bigotimes_\alpha u^{(\alpha)}_{i_\alpha}
    + \sum_\alpha \sum_{i_\alpha} \dot{u}^{(\alpha)}_{i_\alpha}
      \otimes \bU^{(\alpha)}_{i_\alpha}.
\end{equation*}
The Dirac--Frenkel condition~\eqref{eq:DF_tucker_char} then yields
coupled evolution equations for the basis vectors
$u^{(\alpha)}_{i_\alpha}(t) \in V_\alpha$ and the core tensor
$C^{(D)}(t)$, which generalise the MCTDH equations to arbitrary
Tucker ranks in Banach spaces.
\end{remark}

\subsection*{The Hartree manifold}

\begin{example}[Time-dependent Hartree method]
\label{ex:hartree}
For $\fr = \mathbf{1} = (1, \ldots, 1)$, the Tucker manifold
$\Tucker{\mathbf{1}}{\VD}$ consists of elementary tensors
$\bv = \lambda \bigotimes_\alpha v_\alpha$ ($v_\alpha \in V_\alpha$,
$\lambda \in \RR$), called the \emph{Hartree manifold}~\cite{Lubich2008}.
The Dirac--Frenkel ODE~\eqref{eq:DF_tucker} on the Hartree manifold
gives the \emph{time-dependent Hartree equations}. In a Hilbert space
with $\|v_\alpha\|_\alpha = 1$ (unit sphere constraint):
\begin{align*}
    \dot{\lambda}(t) &= \langle \bigotimes_\alpha v_\alpha(t),
        \bF(t, \bv_r(t)) \rangle_D, \\
    \lambda(t)\,\dot{v}_\alpha(t) &= P_{W_\alpha(t)}\bigl[
        (\Id{\alpha} \otimes \langle \cdot, \bigotimes_{\beta \neq \alpha}
        v_\beta(t) \rangle)\bF(t, \bv_r(t))\bigr],
\end{align*}
where $P_{W_\alpha(t)}$ is the orthogonal projection onto the
complement of $\spann\{v_\alpha(t)\}$ in $V_\alpha$. This is
the system solved numerically by MCTDH.
\end{example}

%------------------------------------------------------------
\section{Extension to tree-based formats}
\label{sec:DF:treebased}
%------------------------------------------------------------

The extension of Theorem~\ref{thm:DF_tucker} to $\FTv$ is immediate
given the immersion result of Chapter~\ref{chap:TB_manifold}.

\begin{theorem}[Dirac--Frenkel for tree-based formats,
{\cite[Section~5]{FHN2023}}]
\label{thm:DF_treebased}
Let $\TD$ be a dimension partition tree, $\fr \in \ADTD$ such that
$\FTv$ is a proper set of tree-based tensors. Assume $\VDnorm$ is
reflexive and strictly convex, and that \condINJ\ holds at each
level of $\TD$. Let $\bv_0 \in \FTv$.
\begin{enumerate}
\item[\textup{(Existence)}] There exists a differentiable curve
    $t \mapsto \bv_r(t)$ in $\FTv$ satisfying
    \begin{equation}
    \label{eq:DF_TBT}
        \dot{\bv}_r(t) = \Pmetric{\bv_r(t)}\bigl(\bF(t, \bv_r(t))\bigr),
        \qquad \bv_r(0) = \bv_0,
    \end{equation}
    where $\Pmetric{\bv_r(t)}$ is the metric projection onto
    $T_{\bv_r(t)}\mathfrak{i}(\mathbb{T}_{\bv_r(t)}(\FTv))
    = \bigcap_k \ZD[\PkT{k}{\TD}]{\bv_r(t)}$.
\item[\textup{(Characterisation)}] The curve satisfies~\eqref{eq:DF_TBT}
    if and only if $\dot{\bv}_r(t) \in \bigcap_k \ZD[\PkT{k}{\TD}]{\bv_r(t)}$
    and
    \begin{equation}
    \label{eq:DF_TBT_char}
        \|\dot{\bv}_r(t) - \bF(t, \bv_r(t))\|_D
        = \min_{\dot{\bv} \in \bigcap_k \ZD[\PkT{k}]{\bv_r(t)}}
        \|\dot{\bv} - \bF(t, \bv_r(t))\|_D.
    \end{equation}
\end{enumerate}
\end{theorem}

\begin{proof}
The proof is identical to that of Theorem~\ref{thm:DF_tucker}, with
$\Tucker{\fr}{\VD}$ replaced by $\FTv$ and $\ZD{\bv_r(t)}$ replaced
by $\bigcap_k \ZD[\PkT{k}]{\bv_r(t)}$. The key inputs are:
(i) $\FTv$ is a $\mathcal{C}^\infty$-Banach manifold
(Theorem~\ref{thm:TBT_Banach_Manifold}); (ii) its tangent space is
a split subspace of $\VDnorm$ (Theorem~\ref{thm:TBT_immersion}); and
(iii) the metric projection onto a split subspace in a reflexive
strictly convex space is well-defined and unique.
\end{proof}

\begin{remark}[Algorithmic interpretation]
\label{rem:DF_algorithms}
Theorem~\ref{thm:DF_treebased} provides the mathematical foundation
for a class of algorithms that integrate tensor-format ODEs by evolving
the transfer tensors $\{C^{(\alpha)}\}_{\alpha \in \TD}$ and the leaf
basis vectors $\{u^{(j)}_{i_j}\}_{j \in \Leaves}$ along the
Dirac--Frenkel trajectories. For the TT format, this is the starting
point of the \emph{projector-splitting integrator} of Lubich, Oseledets,
and Vandereycken~\cite{LubichOseledetsVandereycken2015}. For general
tree-based formats, the corresponding algorithm was introduced by
Ceruti, Lubich, and Walach~\cite{CerutiLubichWalach2020}.
\end{remark}

%------------------------------------------------------------
\section{Historical notes}
\label{sec:DF:history}
%------------------------------------------------------------

The Dirac--Frenkel variational principle was proposed independently by
Dirac~\cite{Dirac1930} and Frenkel~\cite{Frenkel1934} in 1930 and 1934
for the time-dependent Schrödinger equation in Hilbert space. The
formulation used here — projection of the vector field onto the tangent
space of a manifold — is that of McLachlan et al.\ and Lubich.

Koch and Lubich~\cite{KochLubich2007} proved existence and uniqueness
of solutions to the Dirac--Frenkel ODE on Tucker manifolds in
finite-dimensional Hilbert spaces, and used it to analyse the MCTDH
method. The extension to infinite-dimensional Banach spaces
(Theorem~\ref{thm:DF_tucker}) was established in~\cite{FHN2019},
where the key difficulty — showing that the tangent space $\ZD{\bv_r(t)}$
is a split subspace of $\VDnorm$, without the Hilbert space structure —
was resolved by the immersion theorem (Theorem~\ref{thm:tucker_immersion}).

The extension to tree-based formats (Theorem~\ref{thm:DF_treebased})
is the main application result of~\cite{FHN2023}: having established
the geometry of $\FTv$ in Chapter~\ref{chap:TB_manifold}, the
Dirac--Frenkel principle follows with no additional effort.

